\documentclass[12pt,a4paper]{article}

\usepackage[T2A]{fontenc}
\usepackage[utf8]{inputenc}
\usepackage[russian]{babel}
\usepackage{amsmath,amssymb,amsthm}
\usepackage{geometry}
\usepackage{array}
\usepackage{longtable}
\usepackage{enumitem}
\usepackage{hyperref}
\usepackage{verbatim}

\geometry{margin=2.2cm}
\setlength{\emergencystretch}{3em}
\setlist{nosep}
\DeclareUnicodeCharacter{03B5}{\ensuremath{\varepsilon}}
\DeclareUnicodeCharacter{2205}{\ensuremath{\varnothing}}

\newtheorem{definition}{Определение}
\newtheorem{remark}{Замечание}

\newcommand{\nullable}{\operatorname{nullable}}
\newcommand{\First}{\operatorname{First}}
\newcommand{\Last}{\operatorname{Last}}
\newcommand{\Follow}{\operatorname{Follow}}
\newcommand{\Pos}{\operatorname{Pos}}
\newcommand{\sym}{\operatorname{sym}}
\newcommand{\move}{\operatorname{move}}

\title{Построение ДКА по регулярному выражению алгоритмом Глушкова}
\author{Проект по курсу ``Теория и реализация языков программирования''}
\date{\today}

\begin{document}

\maketitle

\section{Постановка задачи}

В работе рассматривается учебная реализация преобразования регулярного
выражения в детерминированный конечный автомат. На вход программе подается
строка, задающая регулярное выражение в выбранном синтаксисе. На выходе
формируется LaTeX/TikZ-код диаграммы ДКА. Между этими двумя шагами программа
строит абстрактное синтаксическое дерево, линеаризует литералы, вычисляет
атрибуты Глушкова, строит позиционный автомат и детерминизирует его.

\begin{definition}
Пусть дано регулярное выражение $r$ над конечным алфавитом $\Sigma$.
Требуется построить детерминированный конечный автомат $A_D$ такой, что
$L(A_D)=L(r)$, и вывести код его диаграммы в формате LaTeX/TikZ.
\end{definition}

Проект написан на Haskell и собирается через Cabal. Синтаксический анализ
выполнен с помощью библиотеки Parsec. Реализация не претендует на совместимость
с PCRE или POSIX regular expressions: входной язык специально оставлен
небольшим и формальным, чтобы явно показать этапы алгоритма.

\section{Входной язык и нормализация}

Пользователь вводит выражение в concrete syntax. Приоритет операций принят
обычный: постфиксные операции связывают сильнее конкатенации, а конкатенация
сильнее альтернативы.

\begin{verbatim}
regex   ::= alt
alt     ::= concat ("|" concat)*
concat  ::= postfix+
postfix ::= atom ("*" | "+" | "?")*
atom    ::= literal | escaped | "(" regex ")" | "ε" | "∅"
\end{verbatim}

Пробелы между токенами игнорируются и не входят в алфавит выражения. Обратная
косая черта экранирует служебные символы. Например, \verb|\*| означает
литерал \verb|*|, а не операцию звезды Клини.

Операции \verb|+| и \verb|?| являются синтаксическим сахаром. На этапе разбора
они раскрываются в базовые операции:
\[
  r^+ \equiv rr^*, \qquad r? \equiv \varepsilon \mid r.
\]
После такой нормализации последующие алгоритмы работают уже с ядром регулярных
выражений, а не с полным пользовательским синтаксисом.

\section{Ядро регулярных выражений}

Внутреннее представление использует небольшое ядро:
\[
R ::= \varnothing \mid \varepsilon \mid a \mid R_1R_2
      \mid (R_1 + R_2) \mid R^*.
\]
Здесь $a \in \Sigma$, а знак $+$ в формуле обозначает альтернативу. Во входной
строке для альтернативы используется символ \texttt{|}, а в коде ей
соответствует конструктор \verb|Alt|.

AST выражения задается следующими вариантами:
\begin{itemize}
  \item \verb|EmptySet| --- пустой язык $\varnothing$;
  \item \verb|Epsilon| --- язык, состоящий только из пустого слова;
  \item \verb|Literal Char| --- один символ алфавита;
  \item \verb|Concat [Regex]| --- конкатенация подвыражений;
  \item \verb|Alt [Regex]| --- альтернатива подвыражений;
  \item \verb|Star Regex| --- итерация Клини.
\end{itemize}

\begin{remark}
\verb|Concat| и \verb|Alt| хранятся списками, потому что парсеру удобно
собирать цепочки вроде \verb|abcd| или \texttt{a|b|c} без лишних вложенных узлов.
Семантически эти списки обрабатываются сверткой по бинарным правилам:
$abc$ рассматривается как $((ab)c)$, а $a+b+c$ как $((a+b)+c)$.
Это важно для корректного вычисления \verb|nullable|-подвыражений и отношения
\verb|follow|.
\end{remark}

\section{Линеаризация}

\begin{definition}
Пусть $\Pos(r)=\{1,\ldots,n\}$ --- множество позиций литералов в выражении
$r$. Линеаризация сопоставляет каждому вхождению литерала пару $(p,a)$, где
$p$ --- уникальный номер позиции, а $a$ --- исходный символ.
\end{definition}

Например:
\[
aa \mapsto a_1a_2,
\qquad
a(a \mid a) \mapsto a_1(a_2 \mid a_3).
\]

Линеаризация не является косметическим переименованием. В выражении \verb|aa|
две буквы имеют одинаковый символ, но разные роли: после первой позиции должна
идти вторая, а после второй выражение уже может закончиться. Если различать
только символы, отношение следования невозможно восстановить корректно.

В реализации исходный тип \verb|Regex| не изменяется. Для алгоритма строится
отдельное позиционированное дерево \verb|PositionedRegex|. Таблица
\verb|symbolTable| хранит соответствие \verb|Position -> Char| и позволяет на
поздних этапах вернуться от номеров позиций к символам входного алфавита.

\section{Атрибуты Глушкова}

Для позиционированного выражения вычисляются четыре атрибута.

\begin{definition}
Для регулярного выражения $r$:
\begin{itemize}
  \item $\nullable(r)$ истинно тогда и только тогда, когда
  $\varepsilon \in L(r)$;
  \item $\First(r)$ --- множество позиций, которые могут быть первыми в словах
  языка $L(r)$;
  \item $\Last(r)$ --- множество позиций, которые могут быть последними;
  \item $\Follow_r(i)$ --- множество позиций, которые могут идти сразу после
  позиции $i$ в некотором слове языка $L(r)$.
\end{itemize}
\end{definition}

Для пустого языка:
\[
\nullable(\varnothing)=\mathrm{false},\qquad
\First(\varnothing)=\Last(\varnothing)=\varnothing,\qquad
\Follow=\varnothing.
\]

Для пустого слова:
\[
\nullable(\varepsilon)=\mathrm{true},\qquad
\First(\varepsilon)=\Last(\varepsilon)=\varnothing,\qquad
\Follow=\varnothing.
\]

Для литерала $a_p$ в позиции $p$:
\[
\nullable(a_p)=\mathrm{false},\qquad
\First(a_p)=\Last(a_p)=\{p\},\qquad
\Follow=\varnothing.
\]

Для альтернативы:
\[
\begin{aligned}
\nullable(r+s) &= \nullable(r)\lor\nullable(s),\\
\First(r+s) &= \First(r)\cup\First(s),\\
\Last(r+s) &= \Last(r)\cup\Last(s),\\
\Follow(r+s) &= \Follow(r)\cup\Follow(s).
\end{aligned}
\]

Для конкатенации:
\[
\nullable(rs)=\nullable(r)\land\nullable(s).
\]
\[
\First(rs)=
\begin{cases}
\First(r)\cup\First(s), & \nullable(r),\\
\First(r), & \text{иначе}.
\end{cases}
\]
\[
\Last(rs)=
\begin{cases}
\Last(r)\cup\Last(s), & \nullable(s),\\
\Last(s), & \text{иначе}.
\end{cases}
\]
\[
\Follow(rs)=\Follow(r)\cup\Follow(s)\cup
\{\, i \mapsto \First(s) \mid i\in\Last(r)\,\}.
\]

Для звезды Клини:
\[
\nullable(r^*)=\mathrm{true},
\qquad
\First(r^*)=\First(r),
\qquad
\Last(r^*)=\Last(r).
\]
\[
\Follow(r^*)=\Follow(r)\cup
\{\, i \mapsto \First(r) \mid i\in\Last(r)\,\}.
\]

\begin{remark}
В программе \verb|follow| хранится как \verb|Map Position (Set Position)|.
Это разреженное представление бинарного отношения следования: если для позиции
нет продолжений, отдельная запись в таблице не нужна.
\end{remark}

\section{Позиционный автомат Глушкова}

По вычисленным атрибутам строится позиционный автомат
\[
A_G=(Q,\Sigma,\delta,0,F).
\]
Если после линеаризации выражение имеет $n$ позиций, то
\[
Q=\{0\}\cup\Pos(r).
\]
Состояние $0$ является начальным и не соответствует никакому символу. Каждая
позиция $p\ge 1$ имеет метку $\sym(p)$ из таблицы символов.

Переходы задаются двумя правилами:
\begin{itemize}
  \item если $p\in\First(r)$, то $\delta(0,\sym(p))$ содержит $p$;
  \item если $j\in\Follow_r(i)$, то $\delta(i,\sym(j))$ содержит $j$.
\end{itemize}

Множество принимающих состояний:
\[
F=\Last(r)\cup
\begin{cases}
\{0\}, & \nullable(r),\\
\varnothing, & \text{иначе}.
\end{cases}
\]

Автомат Глушкова не содержит $\varepsilon$-переходов. При этом после замены
позиций обычными символами он может стать недетерминированным.
Например, для выражения \texttt{a|a} из состояния $0$ по символу \verb|a| есть
два перехода: в позицию 1 и в позицию 2.

\section{Детерминизация}

Из позиционного автомата строится ДКА
\[
A_D=(Q_D,\Sigma,\delta_D,q_0,F_D)
\]
методом reachable subset construction. Начальное состояние ДКА соответствует
подмножеству
\[
q_0=\{0\}.
\]
Выбор именно $\{0\}$ отражает состояние позиционного автомата до чтения первого
символа входного слова.

Для $S\subseteq Q$ и символа $c\in\Sigma$:
\[
\move(S,c)=\bigcup_{q\in S}\delta(q,c).
\]
Если $\move(S,c)$ непусто и ранее не встречалось, для него создается новое
состояние ДКА. Таким образом строятся только достижимые подмножества, а не весь
булеан множества $Q$.

Принимающие состояния ДКА определяются правилом:
\[
S\in F_D \iff S\cap F\ne\varnothing.
\]

В реализации ДКА является частичным: отсутствующий переход по очередному
символу означает немедленное отклонение слова. Это не ошибка модели, а
практичное представление автомата для проверки принадлежности. При необходимости
его можно тотализировать добавлением отдельного sink-состояния.

\section{Генерация LaTeX/TikZ}

Генератор выдает standalone-документ LaTeX. Диаграмма строится внутри
окружения \verb|tikzpicture| с использованием стандартных стилей библиотеки
\verb|automata|: \verb|state|, \verb|initial| и \verb|accepting|.

Каждое состояние ДКА $q_i$ выводится как вершина. Перед диаграммой печатаются
комментарии вида:
\begin{verbatim}
% q0 = {0}
% q1 = {1, 3}
\end{verbatim}
Эти строки показывают, как состояния ДКА связаны с подмножествами состояний
позиционного автомата, то есть делают результат subset construction видимым в
сгенерированном файле.

Самопетли и обычные переходы генерируются автоматически. Если несколько
символов ведут из одного состояния в одно и то же состояние, они группируются в
одну подпись ребра. Специальные символы LaTeX в метках экранируются, поэтому
литералы вроде \verb|*|, \verb|_| или \verb|%| можно безопасно выводить в
диаграмме.

\section{Структура реализации}

\begin{longtable}{>{\ttfamily}p{4.6cm}p{9.5cm}}
Модуль & Назначение \\
\hline
Regex.Syntax & AST и pretty-print ядра регулярных выражений. \\
Regex.Parser & Concrete syntax, приоритеты операций и раскрытие сахара. \\
Regex.Linearize & Позиционирование литералов и таблица Position -> Char. \\
Automata.Glushkov & Атрибуты nullable/first/last/follow и позиционный автомат. \\
Automata.DFA & Reachable subset construction, DFA и проверка слова. \\
Render.Tikz & Генерация standalone LaTeX/TikZ-документа. \\
app/Main.hs & CLI-режимы для AST, атрибутов, автоматов, match и TikZ. \\
test/Main.hs & Набор unit/smoke-проверок для основных этапов. \\
\end{longtable}

Такое разделение сохраняет связь между теорией и кодом: каждый математический
этап вынесен в отдельный модуль и может проверяться независимо.

\section{Верификация и тесты}

Текущая версия содержит 69 unit/smoke проверок. Они покрывают:
\begin{itemize}
  \item парсер, приоритеты операций и скобки;
  \item раскрытие синтаксического сахара \verb|+| и \verb|?|;
  \item экранирование служебных символов;
  \item линеаризацию повторяющихся символов;
  \item вычисление \verb|nullable|, \verb|first|, \verb|last| и \verb|follow|;
  \item NFA-переходы для $\varepsilon$, $\varnothing$, \verb|a|, \verb|ab|,
  \verb|a*| и \texttt{a|a};
  \item acceptance-тесты для DFA;
  \item наличие ключевых фрагментов в сгенерированном TikZ.
\end{itemize}

Дополнительно сборка LaTeX проверяет, что генерируемый standalone-файл можно
передать в \verb|pdflatex| без ручного исправления.

\section{Сложность}

Пусть $n$ --- число позиций, то есть число литеральных вхождений в регулярном
выражении. Позиционный автомат Глушкова имеет $n+1$ состояний. Отношение
\verb|follow| в худшем случае может связывать квадратично много пар позиций,
поэтому число переходов позиционного автомата ограничено величиной
\[
O(n^2).
\]

Детерминизация в худшем случае экспоненциальна:
\[
O(2^n)
\]
состояний ДКА. На практике реализация строит только достижимые подмножества,
поэтому для многих выражений фактический автомат оказывается существенно
меньше полного множества подмножеств.

\section{Пример}

Рассмотрим выражение:
\begin{verbatim}
(a|b)*abb
\end{verbatim}

Линеаризация дает позиции:
\begin{verbatim}
1 = 'a'
2 = 'b'
3 = 'a'
4 = 'b'
5 = 'b'
\end{verbatim}

Атрибуты для всего выражения:
\begin{verbatim}
first = {1, 2, 3}
last  = {5}

follow:
1 -> {1, 2, 3}
2 -> {1, 2, 3}
3 -> {4}
4 -> {5}
\end{verbatim}

После детерминизации получаются следующие состояния ДКА:
\begin{verbatim}
q0 = {0}
q1 = {1, 3}
q2 = {2}
q3 = {2, 4}
q4 = {2, 5}
\end{verbatim}

Состояние $q_4$ является принимающим, потому что содержит позицию 5 из
$\Last(r)$. Короткие проверки:
\begin{verbatim}
aabb -> accept
aba  -> reject
\end{verbatim}

\section{Сборка и воспроизведение}

Основные команды для проверки проекта:
\begingroup
\footnotesize
\begin{verbatim}
cabal build
cabal test
cabal run -v0 glushkov-algo -- --match "(a|b)*abb" "aabb"
cabal run -v0 glushkov-algo -- --tikz "(a|b)*abb" > /tmp/dfa.tex
pdflatex -interaction=nonstopmode -halt-on-error -output-directory /tmp /tmp/dfa.tex
\end{verbatim}
\endgroup

Диагностические режимы \verb|--ast|, \verb|--attrs|, \verb|--nfa| и
\verb|--dfa| позволяют посмотреть промежуточные структуры без изменения
основного алгоритма.

\section{Заключение}

В проекте реализован полный основной путь: от строки с регулярным выражением до
ДКА и LaTeX/TikZ-кода диаграммы. Алгоритмическая часть разделена на этапы,
соответствующие построению Глушкова: линеаризация, атрибуты, позиционный
автомат и детерминизация. Минимизация ДКА не выполняется, поскольку она не
требуется постановкой задачи и является отдельным алгоритмом.

\section{Использованные материалы}

\begin{itemize}
  \item Материалы курса ``Теория и реализация языков программирования'',
  репозиторий \url{https://github.com/rbtsv/TRPL}.
  \item Классическая теория регулярных языков и конечных автоматов:
  построение автомата позиций, или алгоритм Глушкова.
  \item Документация Parsec.
  \item Документация Cabal.
  \item Документация TikZ, библиотека \verb|automata|.
\end{itemize}

\end{document}
